\subsection{Conclusiones}

El trabajo evidencia que la Arquitectura Empresarial y la Arquitectura de Seguridad se articulan a través de la gestión de riesgos. El modelado TOGAF (Fases B, C y D) permitió identificar los tres activos críticos del CCI (telemetría, bitácora de decisiones y sistema de despacho) no de forma aislada, sino porque el análisis de negocio, información y tecnología reveló su papel central en la operación. Sin ese mapeo previo no habría sido posible fundamentar las amenazas y su valoración de impacto, de forma que la gestión de riesgos es, en este sentido, el puente que traduce la arquitectura empresarial en decisiones concretas de seguridad.

El caso confirma también la pertinencia del enfoque a sistemas. El propio diseño del CCI (consola unificada, integración por API/bus de eventos, SSO centralizado, principio de ``dato único'') exige tratar la seguridad como una propiedad sostenida en las interacciones entre subsistemas (ITS a bordo, despacho, analítica, información al usuario), y no como controles aislados por aplicación. Los riesgos más críticos se concentran precisamente en esos puntos de integración, por lo que un fallo o manipulación en cualquier subsistema puede propagarse a toda la operación, tal como lo anticipa el principio de continuidad definido en la Fase Preliminar.

Respecto a la valoración de riesgos, destaca que los riesgos priorizados de telemetría, acceso no autorizado al despacho, indisponibilidad del despacho y alteración de la bitácora obtuvieron probabilidad Muy Alta (100~\%), dado el carácter continuo de la operación (24/7, más de 40 millones de registros diarios). Sin embargo, esa probabilidad máxima no derivó en el nivel de riesgo más severo: combinada con impactos Mayor (80~\%) en la mayoría de los casos y Moderado (60~\%) en la bitácora, la severidad resultante los ubicó en el nivel Alto, no Extremo. Esto confirma que, en procesos de operación intensiva y continua, la probabilidad tiende a saturarse, y es el impacto el factor que realmente diferencia y prioriza los riesgos, de ahí que la bitácora (impacto moderado, por no incidir de forma inmediata en la ejecución) se distinga de la telemetría y el despacho (impacto mayor, por su participación directa en las decisiones operativas).

Finalmente, el análisis del impacto de los riesgos sobre los objetivos estratégicos permite definir la estrategia de seguridad en función de estos, la manipulación de la telemetría y la indisponibilidad del despacho amenazan la regulación de flota en tiempo real; el acceso no autorizado compromete la modernización tecnológica y la seguridad de la operación; la alteración de la bitácora afecta la trazabilidad exigida por el PETI institucional. Esta correspondencia legitima la priorización de las fases finales del análisis TOGAF: los quick wins (MFA, registro automático, inventario de interfaces) atienden primero los riesgos de mayor severidad, mientras la redundancia y la auditoría transversal, proyectadas a largo plazo, sostienen la continuidad operativa que es la razón de ser del CCI. La estrategia propuesta, así, no responde a buenas prácticas genéricas, sino que se deriva directamente del análisis de riesgos sobre activos críticos, alineado con la visión definida desde el inicio del ejercicio.